\section{Угрозы и область применимости МКВ}
\label{sec:threats}

\subsection{Угрозы, выявленные в НИРС 1}
% ЧЕРНОВИК: перенести краткий перечень из НИРС 1.
Угрозы систем машинного обучения группируются по этапам конвейера: данные
обучения, модель, данные эксплуатации~\cite{ilyushin2022attacks,papernot2018sok}.

\subsection{Угрозы, которые закрывают МКВ}
\begin{itemize}
    \item раскрытие обучающих данных вычислительным серверам;
    \item раскрытие весов модели при обучении и выводе;
    \item раскрытие входных данных клиента при выводе;
    \item компрометация одного сервера из трёх (при честном большинстве).
\end{itemize}

\subsection{Угрозы, которые МКВ не закрывают}
\begin{itemize}
    \item атаки на результат: определение принадлежности к обучающей
    выборке~\cite{shokri2017mia}, обращение и извлечение модели;
    \item отравление данных владельцем данных, честно выполняющим протокол;
    \item состязательные примеры;
    \item сговор двух серверов из трёх;
    \item побочные каналы, отказ в обслуживании, ошибки реализации.
\end{itemize}

\subsection{Сводное сопоставление}
\begin{table}[H]
    \caption{Сопоставление угроз и возможностей МКВ (заготовка)}
    \label{tab:threat-map}
    \centering
    \begin{tabular}{|p{6cm}|p{3cm}|p{5.5cm}|}
        \hline
        Угроза & Закрывает МКВ & Дополнительная мера \\ \hline
        Раскрытие обучающих данных серверу & да & --- \\ \hline
        Определение принадлежности к выборке & нет & дифференциальная
        приватность, ограничение запросов \\ \hline
        Отравление данных & нет & валидация входных данных \\ \hline
        Сговор большинства серверов & нет & организационное разнесение
        серверов \\ \hline
    \end{tabular}
\end{table}

\subsection{Выводы по разделу}
% ЧЕРНОВИК
